\section{Direct Action 与 Reward Program：系统级比较}

Lecture 的两部分不是互相否定。RT-2 把 action generation 纳入大模型，获得共享表示和快速 inference；Language to Rewards 保留 optimizer，获得 constraint handling 与交互式目标修改。真正的架构选择取决于数据、动作复杂度、控制频率、模型可用性与安全要求。

\subsection{最终总结页：模型、数据与接口}

Final summary 页把第一部分概括为 model consolidation、joint scaling、positive transfer，把第二部分概括为 new interfaces to LLMs。Key takeaway 再强调：rethinking scaling 可以让 robotics 更像 foundation-model training；rethinking interface 可以让 LLM 负责 semantic translation 而不是承担所有低层物理细节。

\lecturefigure{slide-54-final-summary.jpg}{Final summary：Language to Rewards 与 general pattern machines}{官方录像恢复幻灯片 V054；01:09:39}

\lecturefigure{slide-55-key-takeaway.jpg}{Key takeaway：共同扩展模型，也重新设计接口}{官方录像恢复幻灯片 V055；01:10:39}

\subsection{两种路线的 tradeoff table}

下面的表不是选出永久赢家，而是把 component responsibility 写清。系统也可以混合：先用 Language to Rewards 在 simulator 中生成 diverse trajectories，再 distill 到 VLA；或让 VLA 提议 skill、MPC 负责最后一段高精度接触。

\begin{longtable}{@{}L{0.20\textwidth}L{0.32\textwidth}L{0.32\textwidth}@{}}
\toprule
维度 & Direct action / VLA & Reward program + controller \\
\midrule
主要输出 & action tokens 或连续动作参数 & state/action reward code 或 objective \\
语义来源 & web + robot co-fine-tuning & pretrained LLM + robot schema / prompt \\
物理执行 & policy 隐式学 dynamics 与 feedback & optimizer 显式使用 dynamics 与 constraint \\
推理成本 & 一次或少量 autoregressive decoding & 每步需要 rollout / optimization \\
数据需求 & 大量 robot trajectories，支持 distillation & 可零样本写 reward，但需 simulator/model 与 validation \\
优势 & 端到端、快速、可与 VLM 表示共享 & 动作多样、约束明确、可交互修改目标 \\
主要风险 & OOD action、quantization、behavior support、不可解释失败 & reward hacking、code error、model bias、在线优化失败 \\
适合场景 & 高频重复 skill、充足 demonstrations、固定 embodiment & 新语义任务、复杂 constraint、可用 dynamics model \\
\bottomrule
\end{longtable}

\begin{importantbox}{课堂给出的混合方案}
用 Language to Rewards 生成大规模、语义多样的 simulator trajectories，再把这些数据 distill 到 Transformer/VLA，可同时利用 optimizer 的探索能力与 policy 的快速 inference。反方向也可行：VLA 负责高层 skill proposal，MPC 在局部执行并满足碰撞、平衡或力限制。
\end{importantbox}

\teachervoice{Q\&A 中讲者明确比较两条路线：direct action generation 与 autoregressive modeling 高度一致，但可能难以产生训练分布外的 dexterous motion；reward generation 借用了传统 robotics 的 optimization，可加入 safety constraint 并产生更多样动作。这个回答比“端到端或模块化二选一”更接近真实系统设计。}

\subsection{Data 仍是最大的 bottleneck}

即使模型规模增长，机器人领域仍缺少互联网级闭环数据。语言中的 positive transfer 已经普遍到不再被单独强调，而 robotics 仍在寻找早期信号。数据瓶颈不仅是数量：失败、恢复、危险状态、长期任务和不同 embodiment 的 coverage 往往比成功 demonstrations 更稀缺。

\begin{knowledgebox}{数据审计的六个问题}
采集了多少独立 trajectories？失败是否保留？环境和物体是否多样？action distribution 是否覆盖 recovery？不同 operator 是否造成 style bias？evaluation 是否在真正 held-out environment 和 object 上？如果这些问题没有答案，仅报告“总小时数”不足以判断 dataset quality。
\end{knowledgebox}

\subsection{Physical safety 不能只继承 LLM alignment}

最后的学生提问聚焦安全。机器人错误直接作用于人和环境，因此不能只依赖 LLM 拒答。系统需要 hardware torque/force limit、emergency stop、workspace boundary、collision monitor、software interlock、safe controller、simulation test 和 human supervision。讲者还提到团队在探索 constitutional safety，但没有公开细节；讲义只保留这一未来方向，不把未发布系统写成已验证方案。

\begin{warningbox}{“模型不会说危险话”不等于机器人安全}
安全目标至少包括：指令本身是否合法、模型是否理解对象指代、计划是否进入危险区域、低层动作是否超过力/速度限制、传感器失效时是否 fail safe、用户是否能中止。Language alignment 只覆盖其中一部分；hardware 与 controller safety layer 必须独立存在。
\end{warningbox}

\teachervoice{讲者说团队“非常认真”对待 safety，因为物理系统可能伤害人或环境。他举出可脱落 robot finger 作为硬件力限制的例子，也用“把猫放进洗碗机”的歧义指令说明 semantic misunderstanding 会成为 catastrophic physical failure，而不是普通文本 hallucination。}

\subsection{本章小结}

VLA 与 reward-program 路线分别把更多责任交给 learned policy 或 optimizer。前者擅长共享表示与快速 inference，后者擅长显式 constraint 与新任务合成；混合系统往往更现实。无论路线如何，robot data coverage 和独立 safety layers 都不能被 scale 替代。

